\section{מסקנות}
בפרויקט זה פותחה מערכת חיזוי מבוססת למידת מכונה להערכת משתני בית השורשים בחממה, תוך שילוב נתוני אקלים, מיקרו-אקלים, השקיה, דישון, מצב הצמח ומדידות עבר של pH ו-EC.

המסקנות המרכזיות הן:

\begin{itemize}
    \item ניתן לבנות מסד נתונים מאוחד המשלב מקורות מידע הטרוגניים ברזולוציות זמן שונות, ולהפוך אותם לבסיס מתאים לאימון מודל חיזוי.
    \item מודל המיקרו-אקלים מספק תחזית אמינה של תנאי הסביבה בתוך החממה, ומשמש שכבת קלט חשובה למודל בית השורשים.
    \item מודל בית השורשים מצליח לחזות את ערכי pH ו-EC בדיוק טוב, ומשפר את הביצועים ביחס למודל נאיבי המבוסס על שמירת הערך האחרון שנמדד.
    \item חיזוי pH התגלה כיציב יחסית, בעוד שחיזוי EC מאתגר יותר בעיקר בטווחי מוליכות נמוכים.
    \item הגישה ההיברידית, המשלבת \GB{} עם רכיב ליניארי חסין, מאפשרת להתמודד טוב יותר עם שילוב של דפוסים מוכרים ואירועים תפעוליים בעלי השפעה גבוהה.
    \item הממשק הדיגיטלי שפותח מדגים כיצד ניתן להנגיש את המודל למשתמש תפעולי ולהפוך אותו לכלי תומך החלטה ראשוני.
\end{itemize}

בסיכום, הפרויקט מדגים היתכנות לפיתוח \SoftSensor{} לחיזוי pH ו-EC בבית השורשים, ומראה כי למידת מכונה יכולה לתרום לניהול מדויק ופרואקטיבי יותר של השקיה ודישון בחממות.
